\newcommand{\C}{\mathbb{C}}
\newcommand{\R}{\mathbb{R}}
\newcommand{\Z}{\mathbb{Z}}
\newcommand{\Tr}{\operatorname{Tr}}
\newcommand{\Impart}{\operatorname{Im}}
\newcommand{\AF}{\operatorname{AF}}
\newcommand{\AX}{\operatorname{AX}}
\newcommand{\GL}{\operatorname{GL}}

\examyear{2012}
\examera{平成24年}
\examfield{解析}
\examgroup{B}
\examnumber{17}
\examtitle{平成24年 B 第17問}
\examtag{常微分方程式}
\examtag{周期関数}
\examtag{漸近挙動}

次の線形微分方程式を考える。
\[
  (*)\qquad
  \frac{dx(t)}{dt}=\rho(t)x(t)+f(t),\qquad
  x(0)\ne0,\quad t\in[0,+\infty).
\]
ここで，\(\rho(t),f(t)\) は \((-\infty,+\infty)\) で定義された \(\theta>0\) を周期とする連続な周期関数である。
また
\[
  \rho^*=\frac1\theta\int_0^\theta \rho(s)\,ds
\]
と定義する。

\begin{enumerate}
  \item \(\rho^*=0\) である場合，\(f\equiv0\) であれば，解 \(x(t)\) は \(\theta\) を周期とする周期関数であることを示せ。
  \item \(\rho^*=0\) である場合，解 \(x(t)\) に対して，関数 \(y(t)\) を
  \[
    y(t)=e^{-\int_0^t\rho(\zeta)\,d\zeta}x(t)
  \]
  と定義する。このとき，
  \[
    \lim_{t\to+\infty}\frac{y(t)}{t}
      =\frac1\theta\int_0^\theta
      e^{-\int_0^\sigma \rho(\xi)\,d\xi}f(\sigma)\,d\sigma
  \]
  となることを示せ。
  \item \(\rho^*<0\) である場合，関数 \(z(t)\) を以下のように定義する：
  \[
    z(t)=\int_0^{+\infty}
      e^{\rho^*\sigma}\frac{\phi(t)}{\phi(t-\sigma)}f(t-\sigma)\,d\sigma,
  \]
  ただし，
  \[
    \phi(t)=\exp\left(\int_0^t(\rho(\zeta)-\rho^*)\,d\zeta\right)
  \]
  とする。このとき，\(z(t)\) は \(\theta\) を周期とする周期関数で，次を満たすことを示せ。
  \[
    \frac{dz(t)}{dt}=\rho(t)z(t)+f(t),\qquad t\in[0,+\infty).
  \]
  \item \(\rho^*<0\) と仮定する。 \((*)\) の任意の解 \(x(t)\) について，
  \[
    \lim_{t\to+\infty}|x(t)-z(t)|=0
  \]
  となることを示せ。
\end{enumerate}